\documentclass[9pt]{extarticle}
\usepackage[margin=1in, top=0.85in, bottom=0.85in]{geometry}
\usepackage{hyperref}
\usepackage{booktabs}
\usepackage{enumitem}
\usepackage{titlesec}
\usepackage{multicol}
\setlength{\multicolsep}{3pt}
\setlength{\columnsep}{12pt}

\setlength{\parskip}{2pt}
\setlength{\parindent}{0pt}

\titlespacing*{\section}{0pt}{10pt}{5pt}
\titlespacing*{\subsection}{0pt}{8pt}{4pt}

\setlist{noitemsep, topsep=1pt, partopsep=0pt, parsep=0pt}

\begin{document}

\begin{center}
    {\large\textbf{CS 581 --- Nuclear Cybersecurity}} \\[2pt]
    {\normalsize Security for Cyber-Physical \& Nuclear Systems} \\[1pt]
    {\normalsize Boise State University \quad|\quad Fall 2026}
\end{center}
\vspace{2pt}

%------------------------------------------------------------
\section*{Course Information}
%------------------------------------------------------------

\textbf{Location}: Online -- Canvas; \textbf{Time}: Asynchronous \\
Virtual Office Hours: Wed 11am--3pm MT / By Appointment \\
Instructor: Christopher Spirito \\
Email: \href{mailto:cspirito@boisestate.edu}{cspirito@boisestate.edu} ~|~ \href{mailto:chris.spirito@gmail.com}{chris.spirito@gmail.com}

%------------------------------------------------------------
\section*{Course Objectives and Learning Outcomes}
%------------------------------------------------------------

Security is paramount to creating software and systems that operate in
high-consequence environments. This course introduces the fundamentals of
cybersecurity and security engineering through the lens of the nuclear industry
--- one of the most regulated, adversarially targeted, and safety-critical
sectors in the world.

\medskip

Students will move beyond passive study: every module includes a hands-on
\textbf{coding assignment} in which students use an \textbf{AI Coding Assistant}
(e.g., Claude Code, Codex, Antigravity, or equivalent) to design, build,
and iterate on working \textbf{security prototypes}. These prototypes --- threat
models, access control implementations, anomaly detection sketches, compliance
analysis tools, and incident response artifacts --- form a professional portfolio
that students carry out of the course.

\medskip

Topics include vulnerability analysis, common exploits and defenses, applied
cryptography, ICS/SCADA protocol security, nuclear command and control,
insider threat, side-channel attacks, and emerging challenges at the intersection
of AI and nuclear systems. The course consists of pre-recorded lectures,
AI Coding Assistant workshops, individual and group coding activities, a mini-research
experience, a semester project, and practitioner-led Zoom sessions.

\subsection*{Course Topics}

The course follows a deliberate eight-phase progression across 16 weeks,
building from security foundations through to advanced and emerging threats.

\medskip

\begin{multicols}{2}

\textbf{Phase 1 --- Foundations} \textit{(Wks 1--2)}
\begin{itemize}
    \item Computer Security Overview
    \item Psychology of Security
\end{itemize}

\medskip

\textbf{Phase 2 --- Threat Landscape} \textit{(Wks 3--4)}
\begin{itemize}
    \item Opponents \& Adversary Taxonomy
    \item Insider Threat
\end{itemize}

\medskip

\textbf{Phase 3 --- Technical Core} \textit{(Wks 5--7)}
\begin{itemize}
    \item Cryptography \& Protocol Security
    \item Access Control \& Identity
    \item ICS/SCADA \& Digital I\&C
\end{itemize}

\medskip

\textbf{Phase 4 --- Nuclear Context} \textit{(Wks 8--9)}
\begin{itemize}
    \item Physical Protection \& Regulatory Frameworks
    \item Nuclear Command \& Control
    \item Economics of Security
\end{itemize}

\columnbreak

{\setlength{\leftskip}{1.5em}
\textbf{Phase 5 --- Defense \& Detection} \textit{(Wks 11--12)}
\begin{itemize}
    \item Monitoring, Metering \& Anomaly Detection
    \item Supply Chain Security
\end{itemize}

\medskip

\textbf{Phase 6 --- Incident Response} \textit{(Wk 13)}
\begin{itemize}
    \item IR in Nuclear Environments
    \item Zimmerman's 10 Rules for IR
\end{itemize}

\medskip

\textbf{Phase 7 --- Advanced \& Emerging} \textit{(Wks 14--15)}
\begin{itemize}
    \item Side Channels
    \item AI, Autonomy \& SMRs
    \item Emerging Threats
\end{itemize}

\medskip

\textbf{Phase 8 --- Capstone} \textit{(Wk 16)}
\begin{itemize}
    \item Ethics \& Policy
    \item Guest Practitioner Perspectives
\end{itemize}
}

\end{multicols}

%------------------------------------------------------------
\section*{Texts and Materials}
%------------------------------------------------------------

\textbf{Required} \\
\textit{Security Engineering} --- Ross Anderson (2nd and 3rd ed.) \\
Access options (no purchase required):
\begin{itemize}
    \item \textbf{BSU Library (full text):} \href{https://research.ebsco.com/c/fzqa7n/search/details/oetacepj7f?q=security+engineering}{research.ebsco.com} --- log in with your BSU credentials
    \item \textbf{Author's website (selected chapters):} \href{https://www.cl.cam.ac.uk/~rja14/book.html}{cl.cam.ac.uk/\textasciitilde rja14/book.html}
\end{itemize}

\medskip

\textbf{Supplemental} \\
Nuclear-sector primary sources assigned per module, including NRC regulatory
guidance (10 CFR 73.54, RG 5.71), NEI 08-09, IAEA Nuclear Security Series
documents, NIST SP 800-82, and selected academic papers on ICS/SCADA and
cyber-physical system security.

\medskip

\textbf{Required Software} \\
All coding workshops require an AI Coding Assistant with function calling and web
search capability. Obtain access to one before Week 1:
\begin{itemize}
    \item \textbf{Claude Code} (Anthropic) --- \href{https://claude.ai/code}{claude.ai/code}
    \item \textbf{Codex} (OpenAI) --- \href{https://platform.openai.com/}{platform.openai.com}
    \item \textbf{Antigravity} (Google) --- \href{https://antigravity.google/}{antigravity.google}
\end{itemize}
All three are US-based providers with documented data handling policies appropriate
for coursework. A GitHub account is also required; see Module 1 for repository
setup instructions.

%------------------------------------------------------------
\section*{Course Requirements}
%------------------------------------------------------------

\begin{minipage}[t]{0.52\textwidth}
\vspace{0pt}\raggedright
All assessments are real work products.

\medskip

\textbf{Participation} \\
\textit{Attend and engage in at least 8 live Zoom sessions with invited expert practitioners.}

\smallskip
\textbf{Ethics Paper} \\
\textit{Personalized ethical dilemma issued by the instructor in Week 14, drawn from each
student's W1--W8 portfolio artifacts. Students argue a defensible position using primary
sources within a 10-day window.}

\smallskip
\textbf{Workshops \& Coding Assignments} \\
\textit{AI Coding Assistant-driven security prototypes submitted with session logs each module.
Students select a target system from an eleven-system nuclear plant architecture in Week 1
and apply each subsequent workshop to a chosen system; a sequencing plan (graded separately)
ensures all eleven systems appear in the portfolio by Week 14.}

\smallskip
\textbf{Portfolio Capstone} \\
\textit{Integrated executive-level security assessment synthesizing all nine workshop artifacts
into a coherent argument about nuclear plant security posture. Both repos must be in
professional final state by the submission deadline.}

\smallskip
\textbf{Practitioner Engagement} \\
\textit{Written analysis connecting each guest speaker's experience to course material.}

\smallskip
\textbf{Midterm Scenario Exercise} \\
\textit{Personalized nuclear incident scenario at Week 10, tailored to each student's W1--W5
portfolio. Students produce a full incident response package within 72 hours.}
\end{minipage}%
\hfill
\begin{minipage}[t]{0.38\textwidth}
\vspace{0pt}
\begin{tabular}{lr}
\toprule
\textbf{Component} & \textbf{Points} \\
\midrule
System Selection \& Sequencing Plan & 30 \\
Workshops W1--W9 (9 $\times$ 30~pts) & 270 \\
Midterm Scenario Exercise & 100 \\
Ethics Paper & 150 \\
Portfolio Capstone & 250 \\
Participation & 100 \\
Practitioner Engagement & 100 \\
\midrule
\textbf{Total} & \textbf{1000} \\
\bottomrule
\end{tabular}
\end{minipage}

%------------------------------------------------------------
\section*{Grade Policies}
%------------------------------------------------------------

\textbf{Submission} \\
\textit{All assignments must be committed to the student's two course GitHub repositories
(primary: technical artifacts; secondary: role analyses) by 11:59pm MT on the due date.}
We evaluate the last commit before the deadline --- no Canvas file uploads are required for workshops.
Setup instructions are in Module 1. Late submissions are subject to the following
progressive penalty:

\medskip

\begin{minipage}[t]{0.45\textwidth}
\raggedright
\begin{tabular}{ll}
\toprule
\textbf{Lateness} & \textbf{Penalty} \\
\midrule
Up to 24 hours   & 10\% deduction \\
24--48 hours     & 25\% deduction \\
48--72 hours     & 50\% deduction \\
72+ hours        & Not accepted \\
\bottomrule
\end{tabular}
\end{minipage}%
\hfill
\begin{minipage}[t]{0.45\textwidth}
\raggedright
\begin{tabular}{ll}
\toprule
\textbf{Points} & \textbf{Grade} \\
\midrule
$>$899   & A \\
800--899 & B \\
700--799 & C \\
600--699 & D \\
$<$600   & F \\
\bottomrule
\end{tabular}
\end{minipage}

\medskip

Grades are computed as a direct unweighted sum of all course components.
Numerical grades are rounded to the nearest whole number (799.5 becomes 800,
a B; 799.4 remains 799, also a B). Curves, if applied, will only ever move
in your favor. Grade boundaries may be relaxed but never tightened.
Grade summaries are posted on Canvas --- review your scores and report
any discrepancies promptly. Individual assignments may contain optional
extra credit components; no separate extra credit assignments will be offered.

%------------------------------------------------------------
\section*{Warning}
%------------------------------------------------------------

This course covers offensive security techniques with the objective of understanding
how to design defenses. \textbf{Do not apply any technique, tool, or knowledge from
this course against any real system without explicit prior written permission.}
Unauthorized access is illegal, violates BSU network policy, and is incompatible
with the professional standards of this field. When in doubt, stop and ask.

%------------------------------------------------------------
\section*{Course Conduct}
%------------------------------------------------------------

Treat all members of this course --- peers and instructor --- with respect.
Engage with course materials consistently, meet deadlines, and communicate
professionally in all course-related interactions. For email or messages,
include a clear subject line and state your purpose concisely. Active
participation in Zoom sessions with guest practitioners is expected and
contributes directly to your grade. Use professional language in all forums.

%------------------------------------------------------------
\section*{Syllabus Changes}
%------------------------------------------------------------

This syllabus is a guide, not a contract. Deadlines, requirements, and course
structure are subject to change at the instructor's discretion. Changes will
be announced via Canvas and BSU email --- check both regularly.

\end{document}
